\documentclass[../../../include/open-logic-section]{subfiles}

\begin{document}

\olfileid{sth}{z}{story}
\olsection{شرحی دقیق‌تر از روایت}
	
در~\olref[story][approach]{sec} توصیف شوئنفیلد از فرایند تشکیل مجموعه‌ها را
نقل کردیم. اکنون می‌خواهیم چند اصل دیگر را بنویسیم تا این روایت را اندکی
دقیق‌تر سازیم. این اصول از این قرارند:
\begin{enumerate}
\item[] \stageshier. هر مجموعه در مرحله‌ای تشکیل می‌شود.
\item[] \stagesord. مراحل مرتب‌اند: برخی \emph{پیش از} برخی دیگر
می‌آیند.\footnote{در واقع، به‌طور ضمنی فرض خواهیم کرد که مراحل
\emph{خوش‌ترتیب} هستند. معنای این فرض در~\olref[ordinals][]{chap} توضیح
داده می‌شود. این فرضی پرمایه است. در واقع، با استفاده از روشی بسیار
هوشمندانه از~\citet{Scott1974} می‌توان از این فرض \emph{پرهیز کرد} و سپس
آن را \emph{استنتاج کرد}. (این امر همچنین توضیح خواهد داد که چرا باید به
وجود مرحله‌ای آغازین باور داشته باشیم.) اینجا نمی‌توانیم وارد این بحث
شویم؛ برای مطالب بیشتر بنگرید به~\citet{ButtonLT1}.}
\item[] \stagesacc. برای هر مرحلهٔ~$S$ و هر مجموعه‌ای از مجموعه‌هایی که
\emph{پیش از} مرحلهٔ~$S$ تشکیل شده‌اند، در مرحلهٔ~$S$ مجموعه‌ای تشکیل
می‌شود که اعضایش دقیقاً همان مجموعه‌ها هستند. در مرحلهٔ~$S$ هیچ چیز دیگری
تشکیل نمی‌شود.
\end{enumerate}
این‌ها اصولی غیرصوری‌اند، اما می‌توانیم با استفاده از آن‌ها چند اصل موضوع
نظریهٔ مجموعه‌های زرملو را توجیه کنیم.

(باید نکته‌ای احتیاطی را یادآور شویم. هرچند چند اصل موضوع کاملاً معیار را
با نام‌هایی کاملاً معیار عرضه خواهیم کرد، اصول ایتالیکی که اکنون ارائه
کردیم در آثار این حوزه نام مشخصی ندارند. ما صرفاً نام‌هایی بر آن‌ها
نهاده‌ایم که امیدواریم سودمند باشند.)

\end{document}
